\subsection{Graphing \(y = |-2x+8|\)}

Consider the absolute-value function

\[
y = |-2x+8|.
\]

\subsubsection{Finding the Vertex}

The graph bends where the expression inside the
absolute value equals zero:

\[
\begin{aligned}
-2x+8 &= 0 \\
-2x &= -8 \\
x &= 4
\end{aligned}
\]

Substituting \(x=4\):

\[
y = |-2(4)+8| = |0| = 0.
\]

Therefore, the vertex is

\[
\boxed{(4,0)}.
\]

\subsubsection{Creating a Table}

\[
\begin{array}{c|c|c|c}
x & \text{Substitution} & \text{Simplification} & y \\ \hline
2 & |-2(2)+8| & |4| & 4 \\
3 & |-2(3)+8| & |2| & 2 \\
4 & |-2(4)+8| & |0| & 0 \\
5 & |-2(5)+8| & |-2| & 2 \\
6 & |-2(6)+8| & |-4| & 4
\end{array}
\]